\section{Construction of the finite precision tower}
\label{sec:ex-lattice-rg-horizon-geometry}

This section explains the finite-dimensional algebra that turns the selected
quadratic beds and their curvature factor into an exact Wilson tower. It also
records how the beds and towers are produced on the actual schedule. All
matrices act on finite real torus spaces; there is no limiting or spectral
argument in the tower identities.

\subsection{The exact finite-torus block identity}
\label{subsec:ex-lattice-rg-roll}

Let $\mathcal E_f^r$ and $\mathcal E_c^r$ be the real coordinate spaces of
r-form labels on the fine and coarse tori. The first coboundaries are
$D_f: \mathcal E_f^1\to\mathcal E_f^2$ and
$D_c: \mathcal E_c^1\to\mathcal E_c^2$; the next coarse coboundary is
$d_2: \mathcal E_c^2\to\mathcal E_c^3$, with $d_2D_c=0$. Let
$J: \mathcal E_f^2\to\mathcal Y_f$ and
$J': \mathcal E_c^2\to\mathcal Y_c$ be the full precision-factor maps
for the fine and coarse beds, respectively, and write
\[
 W=J^{\mathsf T}J,\qquad W'=J'^{\mathsf T}J',\qquad
 R=W^{-1},\qquad R'=(W')^{-1}.
\]
Assume $W,W'$ are symmetric and the displayed inverses exist. Let $T$ be
the installed two-form transport $T:\mathcal E_c^2\to\mathcal E_f^2$,
and let $A_0:\mathcal E_c^1\to\mathcal E_f^1$ be the selected one-form
transport. The conditional mean $\overline T:\mathcal E_f^1\to\mathcal E_c^1$
and these transports satisfy
\begin{equation}
 T^{\mathsf T}D_f=D_c\overline T,
 \qquad TW'D_c=WD_fA_0,
 \label{eq:ex-lattice-horizon-diagonal}
\end{equation}
where $W$ and $R$ act on $\mathcal E_f^2$, while $W'$ and $R'$
act on $\mathcal E_c^2$. Put $s=\sqrt{k}$, with $k>0$.

Define the curvature remainder by
\begin{equation}
 Z=R'-T^{\mathsf T}RT-k^{-1}D_cD_c^{\mathsf T}.
 \label{eq:ex-lattice-horizon-defect}
\end{equation}
Suppose this remainder has the curvature factorization
\begin{equation}
 W'ZW'=d_2^{\mathsf T}P^{\mathsf T}P d_2.
 \label{eq:ex-lattice-horizon-bleep}
\end{equation}
Let $P:\mathcal E_c^3\to\mathcal E_c^3$ be a real matrix with the
curvature factorization above. Define the coarse projection and its
curvature image by
\[
 \Pi=J'R'J'^{\mathsf T},\qquad
 X=P d_2R'J'^{\mathsf T}.
\]
The three blocks in the tower are
\begin{equation}
 M_F=JRTJ'^{\mathsf T},\qquad
 M_N=s^{-1}D_c^{\mathsf T}J'^{\mathsf T},\qquad
 M_O=\begin{bmatrix}I-\Pi\\X\end{bmatrix}.
 \label{eq:ex-lattice-horizon-blocks}
\end{equation}

\begin{proposition}[Exact block completion]
\label{prop:ex-lattice-horizon-roll}
Assume the inverses and factorization above exist and the two transport
identities \eqref{eq:ex-lattice-horizon-diagonal} hold. Then the blocks in
\eqref{eq:ex-lattice-horizon-blocks} satisfy
\begin{align}
 M_F^{\mathsf T}M_F+M_N^{\mathsf T}M_N+M_O^{\mathsf T}M_O&=I,
 \label{eq:ex-lattice-horizon-gram}\\
 sM_N^{\mathsf T}&=J'D_c,
 \label{eq:ex-lattice-horizon-normal}\\
 M_FJ'D_c&=JD_fA_0,
 \label{eq:ex-lattice-horizon-primal}\\
 M_NJ'D_c&=s^{-1}D_c^{\mathsf T}W'D_c,
 \label{eq:ex-lattice-horizon-kinetic}\\
 M_OJ'D_c&=0,
 \label{eq:ex-lattice-horizon-image}\\
 M_F^{\mathsf T}JD_f+M_N^{\mathsf T}(-s\,\overline T)&=0.
 \label{eq:ex-lattice-horizon-orthogonal}
\end{align}
The three blocks have common domain $\mathcal Y_c$ and respective ranges
$\mathcal Y_f$, $\mathcal E_c^1$, and $\mathcal Y_c\oplus\mathcal E_c^3$.
\end{proposition}

\begin{proof}
The inverse of a symmetric invertible real matrix is symmetric. Indeed,
from $WR=I$ and $W^{\mathsf T}=W$ we obtain $R^{\mathsf T}W=I$, hence
$R^{\mathsf T}=R^{\mathsf T}(WR)=(R^{\mathsf T}W)R=R$; the same argument
applies to $R'$.

The matrix $\Pi$ is the orthogonal projection onto the range of $J'$.
First,
\[
 \Pi^{\mathsf T}=J'R'J'^{\mathsf T}=\Pi.
\]
Second,
\[
 \Pi^2=J'R'(J'^{\mathsf T}J')R'J'^{\mathsf T}
       =J'R'W'R'J'^{\mathsf T}=\Pi,
 \qquad \Pi J'=J'R'W'=J'.
\]

For the curvature block, substitute \eqref{eq:ex-lattice-horizon-bleep}:
\begin{align*}
 X^{\mathsf T}X
 &=J'R'd_2^{\mathsf T}P^{\mathsf T}P d_2R'J'^{\mathsf T}\\
 &=J'R'(W'ZW')R'J'^{\mathsf T}
  =J'ZJ'^{\mathsf T}.
\end{align*}
The two middle factors cancel because $R'W'=W'R'=I$. The other two
Gram terms are
\[
 M_F^{\mathsf T}M_F=J'T^{\mathsf T}RTJ'^{\mathsf T},
 \qquad
 M_N^{\mathsf T}M_N=J'(k^{-1}D_cD_c^{\mathsf T})J'^{\mathsf T}.
\]
By \eqref{eq:ex-lattice-horizon-defect},
$T^{\mathsf T}RT+k^{-1}D_cD_c^{\mathsf T}+Z=R'.$
Consequently the sum of the three block Gram matrices is
\[
 J'R'J'^{\mathsf T}+(I-\Pi)^{\mathsf T}(I-\Pi)
 =\Pi+I-\Pi=I,
\]
using $\Pi^{\mathsf T}=\Pi$ and $\Pi^2=\Pi$. This proves \eqref{eq:ex-lattice-horizon-gram} as an identity on the full
coarse augmented row space $\mathcal Y_c$; no restriction to the range of
$J'$ is made.

The remaining formulas are direct consequences of the block definitions
and \eqref{eq:ex-lattice-horizon-diagonal}. Since $s^2=k$,
$sM_N^{\mathsf T}=J'D_c$. For the fine transport,
\[
 M_FJ'D_c=JRT(J'^{\mathsf T}J')D_c
 =JR(TW'D_c)=JR(WD_fA_0)=JD_fA_0.
\]
For the normal block,
\[
 M_NJ'D_c=s^{-1}D_c^{\mathsf T}W'D_c.
\]
For the image block, its first component is
$(I-\Pi)J'D_c=0$, since $\Pi J'=J'.$ Its second component is
\[
 XJ'D_c=Pd_2R'J'^{\mathsf T}J'D_c
       =Pd_2R'W'D_c=Pd_2D_c=0,
\]
using $R'W'=I$ and $d_2D_c=0$. This proves
\eqref{eq:ex-lattice-horizon-image}.

Finally, the transpose of $M_F$ gives
\[
 M_F^{\mathsf T}JD_f
 =J'T^{\mathsf T}R(J^{\mathsf T}J)D_f
 =J'T^{\mathsf T}D_f.
\]
By $T^{\mathsf T}D_f=D_c\overline T$, this is
$J'D_c\overline T$. Meanwhile
$M_N^{\mathsf T}(s\overline T)=J'D_c\overline T$.
The two terms therefore cancel with the sign in
\eqref{eq:ex-lattice-horizon-orthogonal}. This proves the orthogonality
identity.
\end{proof}

\subsection{The selected schedule, its beds, and the tower estimates}
\label{subsec:ex-lattice-horizon-producer}

The schedule input is a family of priced bed and flight receipts. At each
integer scale $j\geq4$, for every allowed direct precision $p_{\rm direct}
\leq\gamma_j$ and every train of order $j$, the bottom receipt supplies a
bed whose four precision rows are fixed as above, whose five rows have the
stated weighted strip costs, whose floors hold on every torus side, and
whose real precision is sworn positive. It also makes the top bed's
diagonal the identity. For every admissible adjacent pair of orders
$j,t$ (with $t=j$ or $t+1=j$), the flight receipt supplies donated
unmodified primal/two-form transports with the order-$j-2$ difference
bounds, together with a curvature factor of the stated strip cost and
positive lower quadratic bound. The block-length credit requires the
separately selected lower bound on $L$ before its quadratic floor is used.
These receipts are obtained by the one-flight frontier and the uniform bed
pricing; they are construction inputs here, not conclusions of the block
algebra.

Fix the actual level schedule $A$ and a terminal cutoff $N$. Its train at
level $N$ is the constant top train. Descending one level, use the flight
whose wing is the train already constructed at level $i+1$ and whose land
is the train at level $i$. Thus, for every $i<N$, the level-$i$ train is the landing train of the
flight whose wing is the already defined train at level $i+1$. This uses
the actual adjacent block at level $i$; no second schedule is selected.

The bed at each $(N,i)$ is first constructed for the corresponding train.
The source bed proposition fixes its four precision rows to the train's
pulse, dual pulse, amended precision and amended inverse. It constructs the
remaining root row as a square factor of the positive remainder. The
finite-volume root matrix is the vertical stack
\[
 J_p=\begin{bmatrix}\sqrt{p_{\mathrm{direct}}}I\\F\end{bmatrix},
 \qquad J_p^{\mathsf T}J_p
   =p_{\mathrm{direct}}I+F^{\mathsf T}F.
\]
This identity follows by splitting the row index of the product into its
two disjoint blocks. The lower root factor is obtained from the finite
matrix square-factor receipt, so the equality holds at every torus side
and includes the zero modes. The floor receipt then gives symmetry,
inverse identities and the uniform quadratic form bounds required of a
bed.

For a step $i<N$, the flight's donated primal and two-form transports are
the source towers' $A_0$ and $T$. The tangent identities give
\[
 T^{\mathsf T}D_f=D_c\overline T,
 \qquad TW'D_c=WD_fA_0,
\]
where $W,W'$ are the actual adjacent bed metrics and the overline is the
flight's prescribed conditional mean. The flight's rebase receipt supplies
the curvature factorization of the remainder $Z$ in
\eqref{eq:ex-lattice-horizon-defect}. Therefore the preceding block
proposition applies to these actual adjacent beds, not just to an abstract
pair of positive matrices.

The source tower-price theorem also supplies explicit uniform strip costs.
Here is the finite matrix bookkeeping behind that assertion. Let $\alpha$
be the common strip width after narrowing the two bed widths and donated
width; write $\Omega,\Omega'$ for the two bed costs, $V$ for the curvature
factor cost, $L$ for the odd block side, $X$ for the maximum difference
order, $k>0$ for the penalty, and $n_r=|\mathrm{TypeFace}(r)|$ for the
number of $r$-form labels. Put
\begin{align*}
 A&=\max\{\|\sqrt{p_{\rm direct}}\|,\Omega\},
 &B&=\max\{\|\sqrt{q_{\rm direct}}\|,\Omega'\},\\
 C_A&=n_2 A\Omega,& C_B&=n_2B\Omega',&
 C_L&=n_2 C_A\Omega,& C_F&=n_2C_LB,\\
 C_{N}&=k^{-1/2}n_2D_1B,&
 C_{\rm proj}&=n_2C_BB,&
 C_{\rm up}&=n_3VD_2,& C_X&=n_2C_{\rm up}\Omega',\\
 C_O&=\max\{1+C_{\rm proj},n_2C_XB\},&
 M_0&=\max\{\Omega,C_F,C_N,C_O\},
\end{align*}
where
\[
 D_r=\sum_{a\in\mathrm{TypeFace}(r+1)}
       \sum_{b\in\mathrm{TypeFace}(r)}
       \sum_{z\in\mathrm{supp}(\delta_r)}
       e^{\alpha|z|_1}\| (\delta_r)_{ab}(z)\|,
 \qquad r=1,2,
\]
is the finite stencil cost of the $r$th coboundary, where for a cover
displacement $z$ we set $|z|_1=\sum_{j=1}^4|z_j|$. For ordinary cover convolution, the basic weighted row bound can be
checked entrywise. If $C_{ij}=\sum_b A_{ib}*B_{bj}$
is the routed cover convolution and the weighted absolute sums of each
entry of $A$ and $B$ are at most $a_{ib}$ and $b_{bj}$, then
\[
 \sum_z e^{\alpha|z|_1}|C_{ij}(z)|
 \leq \sum_b a_{ib}b_{bj}.
\]
Indeed, write each convolution coefficient as a sum over $x+y=z$, apply
the triangle inequality, use $|x+y|_1\leq|x|_1+|y|_1$, and factor the two
nonnegative sums over $x$ and $y$. Summing the intermediate label $b$
then gives the finite cardinality factor in the strip estimate. Each $n_2$
factor below is this count for a two-form intermediate label; $n_3$ counts
the three-form labels in the curvature factor. A finite coboundary stencil
contributes $D_r$.
The fine-to-coarse composition uses a mixed version of this estimate.
Fix a fine residue $r$ and a fine cover displacement $a$. There is a unique
coarse displacement $\kappa(r,a)$ and a new fine residue
$r'=\operatorname{wrap}(\operatorname{lift}r+a)$ such that
\[
 \operatorname{lift}r+a=\operatorname{lift}r'+L\kappa(r,a).
\]
Coordinatewise, put $u=\operatorname{lift}r$, $v=\operatorname{lift}r'$, and
$n=L$. Both representatives lie in the centered interval, so
$|u-v|\leq n-1$. The displayed identity yields
$n|\kappa_i|\leq |a_i|+n-1$. If $|\kappa_i|>|a_i|$, integrality gives
$|\kappa_i|\geq|a_i|+1$, hence
$n|\kappa_i|\geq n|a_i|+n\geq |a_i|+n$, a contradiction. Thus
$|\kappa(r,a)|_1\leq|a|_1$. In the mixed composition, each $a$ selects
exactly the single residue $r'$ above, and its coarse shift is
$\kappa(r,a)$. Apply the routed-convolution estimate at this fixed residue;
summing the finite intermediate form label costs $n_2$, with no sum over
all possible residues. This is the $C_L=n_2C_A\Omega$ bound, uniform in
$L$.

The fine factor $JR$ has cost $C_A$. Its mixed composition with $T$ has
cost $C_L$, and the final composition with $J'^{\mathsf T}$ has cost
$C_F$; the normal block has cost $C_N$; the projection and curvature rows have
costs $C_{\rm proj},C_X$, respectively. Stacking the projection complement
and curvature rows gives $C_O$. These are precisely the products in the
displayed formulas; the explicit factor $L^{-2}$ in the donated two-form
strip is retained separately. The resulting tower has
\[
 M_N,M_O\text{ strip-cost at most }M_0,
 \qquad M_F\text{ strip-cost at most }M_0L^{-2},
\]
and its $r$th iterated primal and two-form differences have costs at most
$M_0L^{-(1+r)}$ and $M_0L^{-(2+r)}$ for $0\le r\le X$. These latter
bounds are inherited unchanged from the donated tower: the block construction
retains its $A_0,T$ rows, whose input receipts already include every
iterated difference bound through order $X$.

For completeness, the construction has two distinct stages. First, the
schedule supplies a bed for every $(N,i)$ and the recursively defined
flight supplies its one-step data. The bed and hop receipts together with
the exact block proposition above give a horizon with row, floor, crossing,
and roll data. Second, the finite tower-price construction narrows widths
and increases costs, while retaining the original $A_0,T$ rows; the new
rows $M_F,M_N,M_O$ are the block completion just proved. The construction theorem performs this on the actual level object and retains the
same bed family and the actual primal and two-form transports. In
particular its retained-horizon equality is equality of the bed families;
its tower equality preserves the $A_0,T$ rows used by the schedule.

The source contract's remaining numerical receipts are explicit: the
chosen price functions bound the bed strip widths/costs, floor constants,
iterated tower costs, and the one-step crossing and roll identities at all
$N,i$ and finite torus sides. The top bed is the unit diagonal because the
top train is constant. These are exactly the fields of the horizon contract;
they are not consequences of the block identity alone. In the public
construction the level schedule and these receipts are attached to the
preselected run, as stated next.

\subsection{Actual path data and the public binding}
\label{subsec:ex-lattice-horizon-path}

Path data enter through the selected history of that run. Fix a finite
torus, a plaquette face $b$, and a one-form variation $a$. For a path slot
with coefficient $c_{b i}$, cover offset $y_{b i}$, and plaquette label
$p_{b i}$, its physical face is the translate $b+y_{b i}$ modulo the
torus. Define the path linear response by the absolutely convergent sum
$\ell_{\mathrm{path}}(b,a)=\sum_i c_{b i}\,d_1a(b+y_{b i},p_{b i})$. The attachment identity says that, after
summing all path coefficients with fixed $(p,y)$, one obtains the
corresponding cover entry of the root row. The weighted row bound implies
absolute summability of the path coefficients: its exponential weight is
at least one. It also implies absolute summability after multiplying by
the bounded face curl of $a$. Therefore we may first group the path sum
by physical face, and then sum over the finite set of torus faces.
Let $F_W$ be the actual roots block in the selected bed factorization,
acting on periodic two-form fields. Periodizing the cover entries by the
attachment identity gives, for each color component $q$,
\[
 (\ell_{\mathrm{path}}(b,a))_q=\bigl[F_W((d_1a)_q)\bigr](b),
\]
where $F_W$ is the actual root matrix on the torus and $d_1a$ is the
finite-torus coboundary. The right side is a finite matrix-vector
product. This is the entrywise path-to-root identity used in the physical
quadratic-form computation; the absolute bound justifies the only
infinite regrouping.

For the selected RG history, the construction receipt gives the attachment
identity for every cutoff $N$, level $i\leq N$, and period label $p$,
between the bed $H_{N,i}$ and the path at that same history, cutoff, level,
and period. The public adapter carries the original levels and retained
horizon into the component theorem; it does not choose another schedule.
Thus the path-to-root identity above applies to the same bed rows that
enter the tower calculation. The weighted row bound is an additional
priced path input: it is what supplies absolute summability, and is not
concluded from attachment alone.

The matrix identities proved here discharge the tower side of the
selected construction. The path attachment and absolute bounds are
separate receipts of the selected path history, and are kept explicit in
the source crosswalk rather than inferred from the matrix bounds alone.
